\section{Conclusion}
\label{sec:conclusion}

Single sign-on gave students one login for many services, and with it a default in which every service receives the same identifier and the same personal claims, with no stated purpose and no record. Starting from the survey of privacy strategies for identity management by Werner, Westphall and Westphall, we found that accountability and use limitation are the least covered of its nine characteristics, and that its challenge of ensuring that privacy policies are implemented remained open.

TrustTrail addresses these gaps for an OpenID Connect federation. Keycloak caps which disclosure profile each service can obtain and issues pairwise pseudonyms without a shared session identifier; TrustTrail releases attribute values only under a purpose and deletion date that the service has signed, records every release and every refused rule break in a hash-chained ledger, anchors the ledger's Merkle roots on Ethereum, and lowers the tier of services that break the rules. In our evaluation, the privacy layer added a median of 3.6\,ms to a login, all six tamper attacks were detected by at least one check, all 45 scripted violations were detected with no false positives, and each anchor after the first cost 77{,}268 gas regardless of batch size.

TrustTrail does not make a service delete data, and it does not hide logins from the identity provider. It makes the promises of services explicit, checks them at every release, and keeps a record that neither the services nor the operator can change without detection once it is anchored.
